\documentclass[a4paper,12pt]{article}
\usepackage[utf8]{inputenc}
\usepackage[T1]{fontenc}
\usepackage{amsmath, amssymb}
\usepackage{xcolor}
\usepackage[most]{tcolorbox}
\usepackage{geometry}
\usepackage[hidelinks]{hyperref}
\usepackage[turkish]{babel}

\geometry{margin=0.7in}

% Turkish babel fix for active characters in key-value pairs
\shorthandoff{=}
\shorthandoff{:}

\definecolor{cardblue}{RGB}{41, 128, 185} % Modern Blue
\definecolor{cardgray}{RGB}{248, 249, 250}
\definecolor{sorucolor}{RGB}{44, 62, 80}
\definecolor{cevapcolor}{RGB}{39, 174, 96}

\newtcolorbox{learningcard}{
    enhanced,
    colback=white,
    colframe=cardblue!70!white, % Thin blue border like the image
    boxrule=0.6pt,
    sharp corners, % Sharp corners as seen in the image
    left=15pt,
    right=15pt,
    top=12pt,
    bottom=12pt,
    middle=15pt,
    before skip=2em,
    after skip=2em,
    segmentation style={cardblue!10!white, line width=0.5pt}, % Subtle separator
    shadow={2pt}{-2pt}{0pt}{black!5!white}
}

% Labels matching the user's image (Red "Soru", Green "Cevap")
% Numaralandırma sistemi
\newcounter{soruno}
\newcommand{\SoruHeader}{\stepcounter{soruno}{\textbf{\color{red!80!black}Soru \thesoruno:}}\ }
\newcommand{\CevapHeader}{{\textbf{\color{green!60!black}Cevap:}}\ }


\begin{document}

\title{\Huge \textbf{Makine Öğrenmesi} \\ \Large Profesyonel Öğrenme Kartları}
\author{Oğuz Taşdemir}
\date{\today}

\maketitle
\vfill

\begin{center}
    \large
    Bu doküman, teknik mülakat hazırlığı ve akademik tekrar için tasarlanmış \\ \textbf{121 adet} çift taraflı (soru-cevap) öğrenme kartını içerir.
\end{center}

\newpage
\tableofcontents
\newpage

\section{Öğrenme Kartları}

\begin{learningcard}
\SoruHeader
$R^2$ (R-Kare) skoru neyi ölçer? \vspace{10pt}
\tcblower
\CevapHeader
Bağımsız değişkenlerin, bağımlı değişkendeki varyansın ne kadarını açıkladığını.
\end{learningcard}

\begin{learningcard}
\SoruHeader
Adam optimizer'ın temel karakteristiği nedir?
\par\vspace{\baselineskip} \tcblower
\CevapHeader
Adaptif öğrenme hızı kullanarak her parametre için güncellemeyi ayarlar.
\end{learningcard}

\begin{learningcard}
\SoruHeader
Açıklanan Varyans Oranı (EVR) formülü nedir?
\par\vspace{\baselineskip} \tcblower
\CevapHeader
$EVR(k) = \frac{\sum_{i=1}^k \lambda_i}{\sum_{i=1}^d \lambda_i}$
\end{learningcard}

\begin{learningcard}
\SoruHeader
Ağırlıkların hata payına göre güncellenmesini sağlayan matematiksel kural nedir?
\par\vspace{\baselineskip} \tcblower
\CevapHeader
Zincir Kuralı (Chain Rule) tabanlı Geri Yayılım.
\end{learningcard}

\begin{learningcard}
\SoruHeader
Batch Normalization'ın temel amacı nedir?
\par\vspace{\baselineskip} \tcblower
\CevapHeader
Eğitimi hızlandırmak ve gradyan sorunlarını azaltmak için girdi değerlerini normalize etmek.
\end{learningcard}

\begin{learningcard}
\SoruHeader
Bias (Önyargı) ve Varyans arasındaki ödünleşim (trade-off) neyi ifade eder?
\par\vspace{\baselineskip} \tcblower
\CevapHeader
Model karmaşıklığı arttıkça biasın azalması ancak varyansın artması durumunu.
\end{learningcard}

\begin{learningcard}
\SoruHeader
Bias-Variance ödünleşimi açısından 'Underfitting' (Eksik Öğrenme) ne zaman gerçekleşir?
\par\vspace{\baselineskip} \tcblower
\CevapHeader
Model çok basit olduğunda ve verideki temel yapıyı öğrenemediğinde (Yüksek Bias).
\end{learningcard}

\begin{learningcard}
\SoruHeader
Bilgi Kazancı (Information Gain) hesaplanırken hangi kavram temel alınır?
\par\vspace{\baselineskip} \tcblower
\CevapHeader
Entropi (H).
\end{learningcard}

\begin{learningcard}
\SoruHeader
Boyut düşürme neden 'Gürültü Azaltma' (Noise Reduction) sağlar?
\par\vspace{\baselineskip} \tcblower
\CevapHeader
Az varyans taşıyan bileşenler (genellikle gürültü) atıldığı için sinyal/gürültü oranı artar.
\end{learningcard}

\begin{learningcard}
\SoruHeader
Boyut düşürme yöntemlerinden PCA (Temel Bileşen Analizi) hangi tür öğrenme kategorisine girer?
\par\vspace{\baselineskip} \tcblower
\CevapHeader
Denetimsiz (Unsupervised) Öğrenme.
\end{learningcard}

\begin{learningcard}
\SoruHeader
Boyut düşürmede 'Kernel PCA' hangi tür veriler için uygundur?
\par\vspace{\baselineskip} \tcblower
\CevapHeader
Doğrusal olarak ayrılamayan, karmaşık yapılara sahip veriler için.
\end{learningcard}

\begin{learningcard}
\SoruHeader
Boyut düşürmede 'Özellik Seçme' (Feature Selection) ile 'Özellik Çıkarma' arasındaki temel fark nedir?
\par\vspace{\baselineskip} \tcblower
\CevapHeader
Özellik seçme orijinal değişkenlerin bir alt kümesini korurken, özellik çıkarma veriyi yeni yapay bileşenlere dönüştürür.
\end{learningcard}

\begin{learningcard}
\SoruHeader
Boyutların Laneti (Curse of Dimensionality) nedeniyle yüksek boyutlarda veri noktaları arasındaki mesafe nasıl değişir?
\par\vspace{\baselineskip} \tcblower
\CevapHeader
Tüm noktalar birbirine neredeyse eşit uzaklıkta görünür ve 'yakınlık' kavramı anlamını yitirir.
\end{learningcard}

\begin{learningcard}
\SoruHeader
Boyutların Laneti (Curse of Dimensionality) veri noktaları arasındaki mesafeyi nasıl etkiler?
\par\vspace{\baselineskip} \tcblower
\CevapHeader
Boyut arttıkça noktalar arası mesafe birbirine çok benzer hale gelir ve en yakın komşu kavramı anlamını yitirir.
\end{learningcard}

\begin{learningcard}
\SoruHeader
DBSCAN algoritması geleneksel merkez tabanlı algoritmalardan hangi özelliği ile ayrılır?
\par\vspace{\baselineskip} \tcblower
\CevapHeader
Yoğunluk farklarına dayanarak karmaşık şekilli kümeleri bulabilmesi.
\end{learningcard}

\begin{learningcard}
\SoruHeader
DBSCAN gibi yoğunluk tabanlı kümeleme algoritmalarının klasik yöntemlere göre en büyük avantajı nedir?
\par\vspace{\baselineskip} \tcblower
\CevapHeader
Düzensiz ve karmaşık şekilli kümeleri tespit edebilmeleri ve gürültüyü (aykırı değerleri) ayırt edebilmeleridir.
\end{learningcard}

\begin{learningcard}
\SoruHeader
Destek Vektör Makineleri (SVM) algoritmasının temel optimizasyon hedefi nedir?
\par\vspace{\baselineskip} \tcblower
\CevapHeader
Sınıflar arasındaki boşluğu (margin) maksimize eden karar hiperdüzlemini bulmak.
\end{learningcard}

\begin{learningcard}
\SoruHeader
Destek Vektör Makineleri'nde (SVM) 'Margin' neyi ifade eder?
\par\vspace{\baselineskip} \tcblower
\CevapHeader
Karar hiperdüzlemi ile en yakın veri noktaları arasındaki mesafe.
\end{learningcard}

\begin{learningcard}
\SoruHeader
Doğrusal Diskriminant Analizi'nin (LDA) temel amacı nedir?
\par\vspace{\baselineskip} \tcblower
\CevapHeader
Sınıflar arası ayrımı (separability) maksimize etmek.
\end{learningcard}

\begin{learningcard}
\SoruHeader
Doğrusal olmayan yapıdaki veriler için PCA'nın hangi türevi 'Kernel Trick' kullanarak çalışır?
\par\vspace{\baselineskip} \tcblower
\CevapHeader
Kernel PCA (Çekirdek Tabanlı PCA).
\end{learningcard}

\begin{learningcard}
\SoruHeader
Doğrusal regresyon veya lojistik regresyonun çözemediği klasik mantıksal problem hangisidir?
\par\vspace{\baselineskip} \tcblower
\CevapHeader
XOR Paradoksu.
\end{learningcard}

\begin{learningcard}
\SoruHeader
Düzeltilmiş R-Kare (Adjusted $R^2$) neden standart $R^2$'den daha güvenilirdir?
\par\vspace{\baselineskip} \tcblower
\CevapHeader
Modele eklenen alakasız değişkenler için penalizasyon uygulayarak 'yalancı başarıyı' engellediği için.
\end{learningcard}

\begin{learningcard}
\SoruHeader
EM (Expectation-Maximization) algoritması kümeleri nasıl tanımlar?
\par\vspace{\baselineskip} \tcblower
\CevapHeader
Veriyi farklı olasılık dağılımlarının (genellikle Gauss) bir bileşimi olarak varsayar ve 'yumuşak' (soft) kümeleme yapar.
\end{learningcard}

\begin{learningcard}
\SoruHeader
Evrensel Yaklaşım Teoremi (Universal Approximation Theorem) neyi savunur?
\par\vspace{\baselineskip} \tcblower
\CevapHeader
Yeterince geniş tek bir gizli katmana sahip MLP'nin her türlü sürekli fonksiyonu temsil edebileceğini.
\end{learningcard}

\begin{learningcard}
\SoruHeader
Geri Yayılım (Backpropagation) algoritmasının temel amacı nedir?
\par\vspace{\baselineskip} \tcblower
\CevapHeader
Çıkış katmanındaki hatayı zincir kuralı kullanarak geriye doğru dağıtmak ve ağırlıkları güncellemek.
\end{learningcard}

\begin{learningcard}
\SoruHeader
Gini Safsızlığı formülü nedir?
\par\vspace{\baselineskip} \tcblower
\CevapHeader
$Gini(t) = 1 - \sum_k p_k^2$
\end{learningcard}

\begin{learningcard}
\SoruHeader
Gradient Boosting yönteminin Random Forest'tan temel farkı nedir?
\par\vspace{\baselineskip} \tcblower
\CevapHeader
Ağaçların paralel değil, bir önceki ağacın hatalarını (artıkları) öğrenecek şekilde ardışıl olarak eğitilmesi.
\end{learningcard}

\begin{learningcard}
\SoruHeader
Gradient Descent algoritmasındaki 'Öğrenme Oranı' ($\alpha$) neyi belirler?
\par\vspace{\baselineskip} \tcblower
\CevapHeader
Minimuma ulaşmak için her adımda atılan güncellemelerin büyüklüğünü.
\end{learningcard}

\begin{learningcard}
\SoruHeader
Gradyan İnişi (Gradient Descent) algoritmasında öğrenme hızı ($\alpha$) neyi belirler?
\par\vspace{\baselineskip} \tcblower
\CevapHeader
Parametre güncellemeleri sırasında atılan adımların büyüklüğünü.
\end{learningcard}

\begin{learningcard}
\SoruHeader
Hiyerarşik kümelemede 'Complete Linkage' mesafeyi nasıl ölçer?
\par\vspace{\baselineskip} \tcblower
\CevapHeader
İki kümenin elemanları arasındaki en büyük mesafeyi temel alarak.
\end{learningcard}

\begin{learningcard}
\SoruHeader
Hiyerarşik kümelemede 'Dendrogram' neyi temsil eder?
\par\vspace{\baselineskip} \tcblower
\CevapHeader
Veri noktaları arasındaki birleşme hiyerarşisini ve kümeler arasındaki mesafe ilişkilerini.
\end{learningcard}

\begin{learningcard}
\SoruHeader
Hiyerarşik kümelemede 'Single Linkage' (Tekil Bağlantı) mesafeyi nasıl tanımlar?
\par\vspace{\baselineskip} \tcblower
\CevapHeader
İki küme arasındaki en yakın eleman çiftleri arasındaki mesafe.
\end{learningcard}

\begin{learningcard}
\SoruHeader
Hiyerarşik kümelemede 'Single Linkage' (Tekli Bağlantı) kuralı nedir?
\par\vspace{\baselineskip} \tcblower
\CevapHeader
İki küme arasındaki mesafeyi, kümelerin birbirine en yakın eleman çifti arasındaki mesafe olarak tanımlar.
\end{learningcard}

\begin{learningcard}
\SoruHeader
İki nokta arasındaki en kısa doğrusal mesafeyi ölçen metriğe ne denir?
\par\vspace{\baselineskip} \tcblower
\CevapHeader
Öklid (Euclidean) Mesafesi.
\end{learningcard}

\begin{learningcard}
\SoruHeader
İkili sınıflandırma (Binary Classification) için kullanılan temel kayıp fonksiyonu (loss function) nedir?
\par\vspace{\baselineskip} \tcblower
\CevapHeader
Binary Cross-Entropy Loss.
\end{learningcard}

\begin{learningcard}
\SoruHeader
k-En Yakın Komşu (k-NN) algoritması neden 'Tembel Öğrenme' (Lazy Learning) olarak adlandırılır?
\par\vspace{\baselineskip} \tcblower
\CevapHeader
Eğitim aşamasında bir model kurmayıp, tüm hesaplamayı tahmin anında yaptığı için.
\end{learningcard}

\begin{learningcard}
\SoruHeader
K-Katlı Çapraz Doğrulama (K-Fold Cross Validation) yönteminin amacı nedir?
\par\vspace{\baselineskip} \tcblower
\CevapHeader
Veriyi farklı parçalara bölerek modelin genelleme yeteneğini daha güvenilir bir şekilde ölçmek.
\end{learningcard}

\begin{learningcard}
\SoruHeader
k-NN'de $K$ sayısının çok küçük seçilmesi hangi soruna yol açar?
\par\vspace{\baselineskip} \tcblower
\CevapHeader
Yüksek varyans ve aşırı öğrenme (overfitting).
\end{learningcard}

\begin{learningcard}
\SoruHeader
K-Ortalamalar (K-Means) algoritmasında 'Inertia' metriği neyi ifade eder?
\par\vspace{\baselineskip} \tcblower
\CevapHeader
Veri noktalarının kendi küme merkezlerine olan uzaklıklarının kareleri toplamını.
\end{learningcard}

\begin{learningcard}
\SoruHeader
K-Ortalamalar (K-Means) algoritmasının optimizasyon hedefi nedir?
\par\vspace{\baselineskip} \tcblower
\CevapHeader
Küme içi varyasyonu (Inertia) minimize etmek.
\end{learningcard}

\begin{learningcard}
\SoruHeader
Karar ağaçlarında 'Dallanma Kriteri' olarak ne seçilir?
\par\vspace{\baselineskip} \tcblower
\CevapHeader
En yüksek bilgi kazancını (Information Gain) sağlayan özellik ve eşik değeri.
\end{learningcard}

\begin{learningcard}
\SoruHeader
Karar Ağaçlarında 'Gini Safsızlığı' (Gini Impurity) neyi ölçer?
\par\vspace{\baselineskip} \tcblower
\CevapHeader
Bir düğümdeki verilerin sınıfsal homojenliğini (karışıklığını).
\end{learningcard}

\begin{learningcard}
\SoruHeader
Karar ağaçlarında dallanma kriteri olarak kullanılan saflık ölçüsü nedir?
\par\vspace{\baselineskip} \tcblower
\CevapHeader
Gini Safsızlığı veya Entropi.
\end{learningcard}

\begin{learningcard}
\SoruHeader
Karmaşıklık Matrisinde (Confusion Matrix) 'F1 Skoru' nasıl hesaplanır?
\par\vspace{\baselineskip} \tcblower
\CevapHeader
Hassasiyet (Precision) ve Duyarlılığın (Recall) harmonik ortalaması alınarak.
\end{learningcard}

\begin{learningcard}
\SoruHeader
KD-Ağaçları (KD-Trees) uzayı bölmek için hangi geometrik yapıyı kullanır?
\par\vspace{\baselineskip} \tcblower
\CevapHeader
Eksenlere dik hiperdüzlemler.
\end{learningcard}

\begin{learningcard}
\SoruHeader
KD-Ağaçları hangi boyut aralığında verimli arama yapabilmektedir?
\par\vspace{\baselineskip} \tcblower
\CevapHeader
Genellikle 2 ile 8 boyut arasında.
\end{learningcard}

\begin{learningcard}
\SoruHeader
KNN algoritmasında 'K' değerinin çok küçük seçilmesi hangi soruna yol açar?
\par\vspace{\baselineskip} \tcblower
\CevapHeader
Aşırı öğrenme (Overfitting/Yüksek Varyans).
\end{learningcard}

\begin{learningcard}
\SoruHeader
Kök Ortalama Kare Hata (RMSE) metriğinin MAE'ye göre yorumlama avantajı nedir?
\par\vspace{\baselineskip} \tcblower
\CevapHeader
Hata birimini orijinal hedef değişken birimine (Örn: Dolar, Derece) geri döndürmesi.
\end{learningcard}

\begin{learningcard}
\SoruHeader
Kümeleme (Clustering) analizinin temel amacı nedir?
\par\vspace{\baselineskip} \tcblower
\CevapHeader
Grup içi benzerliği en yüksek, gruplar arası farklılığı en fazla olacak şekilde veriyi alt gruplara ayırmak.
\end{learningcard}

\begin{learningcard}
\SoruHeader
Kümeleme Analizi (Clustering) hangi tür öğrenme yöntemidir?
\par\vspace{\baselineskip} \tcblower
\CevapHeader
Gözetimsiz (Unsupervised) Öğrenme.
\end{learningcard}

\begin{learningcard}
\SoruHeader
Kümeleme analizinde 'Öklid Mesafesi' nasıl tanımlanır?
\par\vspace{\baselineskip} \tcblower
\CevapHeader
İki nokta arasındaki en kısa, doğrusal mesafe.
\end{learningcard}

\begin{learningcard}
\SoruHeader
Kümeleme sonuçlarını ağaç benzeri bir yapıda gösteren diyagrama ne ad verilir?
\par\vspace{\baselineskip} \tcblower
\CevapHeader
Dendrogram.
\end{learningcard}

\begin{learningcard}
\SoruHeader
Kümelemede 'Uzman Doğrulaması' (Sanity Check) neden gereklidir?
\par\vspace{\baselineskip} \tcblower
\CevapHeader
İstatistiksel olarak bulunan kümelerin gerçek dünya iş mantığına (Örn: Pazar segmentasyonu) uygunluğunu denetlemek için.
\end{learningcard}

\begin{learningcard}
\SoruHeader
L1 Düzenlileştirmesi (Lasso) modelde nasıl bir etki yaratır?
\par\vspace{\baselineskip} \tcblower
\CevapHeader
Bazı katsayıları tam olarak sıfıra eşitleyerek özellik seçimi (feature selection) yapar.
\end{learningcard}

\begin{learningcard}
\SoruHeader
L2 Düzenlileştirmesi (Ridge) maliyet fonksiyonuna ne ekler?
\par\vspace{\baselineskip} \tcblower
\CevapHeader
Ağırlıkların karelerinin toplamını ($\lambda \cdot \sum w_j^2$).
\end{learningcard}

\begin{learningcard}
\SoruHeader
LDA (Doğrusal Diskriminant Analizi) ile PCA arasındaki en temel kavramsal fark nedir?
\par\vspace{\baselineskip} \tcblower
\CevapHeader
PCA denetimsiz (etiketsiz) bir yöntemken, LDA sınıfları birbirinden ayırmak için etiketli veri kullanan denetimli bir yöntemdir.
\end{learningcard}

\begin{learningcard}
\SoruHeader
LDA'da 'Sınıf İçi Saçılım' ($S_W$) matrisinin minimize edilmesi ne anlama gelir?
\par\vspace{\baselineskip} \tcblower
\CevapHeader
Aynı sınıfa ait örneklerin birbirine olabildiğince yakınlaşması.
\end{learningcard}

\begin{learningcard}
\SoruHeader
LDA'da maksimum bileşen sayısı ne ile sınırlıdır?
\par\vspace{\baselineskip} \tcblower
\CevapHeader
Sınıf sayısı $K$ olmak üzere, en fazla $K-1$ bileşen üretilebilir.
\end{learningcard}

\begin{learningcard}
\SoruHeader
LDA'da üretilebilecek maksimum bileşen sayısı ne ile sınırlıdır?
\par\vspace{\baselineskip} \tcblower
\CevapHeader
Sınıf sayısı eksi bir ($K-1$).
\end{learningcard}

\begin{learningcard}
\SoruHeader
Lojistik Regresyon modelinde çıktı olasılığını $[0, 1]$ aralığına sıkıştıran fonksiyon hangisidir?
\par\vspace{\baselineskip} \tcblower
\CevapHeader
Sigmoid (Lojistik) fonksiyonu: $\sigma(z) = \frac{1}{1 + e^{-z}}$.
\end{learningcard}

\begin{learningcard}
\SoruHeader
Lojistik regresyonda katsayıların aşırı büyümesini engellemek için kullanılan L2 düzenlileştirmesinin diğer adı nedir?
\par\vspace{\baselineskip} \tcblower
\CevapHeader
Ridge Düzenlileştirme.
\end{learningcard}

\begin{learningcard}
\SoruHeader
Manhattan (City-block) mesafesi hangi kurala göre hesaplanır?
\par\vspace{\baselineskip} \tcblower
\CevapHeader
Koordinat eksenleri boyunca hareket edilerek hesaplanan mutlak farkların toplamı.
\end{learningcard}

\begin{learningcard}
\SoruHeader
Manhattan mesafesi nasıl hesaplanır?
\par\vspace{\baselineskip} \tcblower
\CevapHeader
Koordinat eksenleri boyunca mutlak farkların toplamı ile.
\end{learningcard}

\begin{learningcard}
\SoruHeader
MAPE metriği hangi durumlarda hatalı veya tanımsız sonuçlar üretir?
\par\vspace{\baselineskip} \tcblower
\CevapHeader
Gerçek değerler sıfır veya sıfıra çok yakın olduğunda.
\end{learningcard}

\begin{learningcard}
\SoruHeader
Mevcut özelliklerden bir alt küme seçerek orijinal değişkenlerin korunduğu işleme ne denir?
\par\vspace{\baselineskip} \tcblower
\CevapHeader
Özellik Seçme (Feature Selection).
\end{learningcard}

\begin{learningcard}
\SoruHeader
Modelin eğitim verisine aşırı uyum sağlaması sonucu test verisinde düşük performans göstermesi durumuna ne denir?
\par\vspace{\baselineskip} \tcblower
\CevapHeader
Overfitting (Aşırı Öğrenme).
\end{learningcard}

\begin{learningcard}
\SoruHeader
Naive Bayes algoritmasının 'Naive' (Saf) olarak adlandırılmasının sebebi nedir?
\par\vspace{\baselineskip} \tcblower
\CevapHeader
Özelliklerin sınıf verildiğinde birbirinden tamamen bağımsız olduğunu varsaymasıdır.
\end{learningcard}

\begin{learningcard}
\SoruHeader
Naive Bayes türlerinden hangisi sürekli (continuous) değişkenler için uygundur?
\par\vspace{\baselineskip} \tcblower
\CevapHeader
Gaussian Naive Bayes.
\end{learningcard}

\begin{learningcard}
\SoruHeader
Neden 'Düzeltilmiş R-Kare' (Adjusted $R^2$) metriği standart R-Kare'den daha güvenilirdir?
\par\vspace{\baselineskip} \tcblower
\CevapHeader
Modele eklenen anlamsız değişkenler için cezalandırma yaparak sahte başarıyı engellediği için.
\end{learningcard}

\begin{learningcard}
\SoruHeader
Neden PCA uygulanmadan önce verinin standartlaştırılması (z-score scaling) önerilir?
\par\vspace{\baselineskip} \tcblower
\CevapHeader
Ölçeği büyük olan değişkenlerin varyans hesaplamasını domine etmesini engellemek için.
\end{learningcard}

\begin{learningcard}
\SoruHeader
Ortalama Kare Hata (MSE) metriği hataları nasıl cezalandırır?
\par\vspace{\baselineskip} \tcblower
\CevapHeader
Hataların karesini aldığı için büyük sapmaları çok daha sert (üstel olarak) cezalandırır.
\end{learningcard}

\begin{learningcard}
\SoruHeader
Ortalama Kare Hata (MSE) metriğinin MAE'den temel farkı nedir?
\par\vspace{\baselineskip} \tcblower
\CevapHeader
Hataların karesini aldığı için büyük hataları daha sert cezalandırması.
\end{learningcard}

\begin{learningcard}
\SoruHeader
Ortalama Mutlak Hata (MAE) metriğinin en güçlü yönü nedir?
\par\vspace{\baselineskip} \tcblower
\CevapHeader
Aykırı değerlere (outliers) karşı dirençli ve stabil bir ölçüm sunması.
\end{learningcard}

\begin{learningcard}
\SoruHeader
Ortalama Mutlak Yüzde Hata (MAPE) metriğinin temel kısıtlaması nedir?
\par\vspace{\baselineskip} \tcblower
\CevapHeader
Gerçek değerler sıfıra çok yakın olduğunda sonucun matematiksel olarak belirsizleşmesi veya patlaması.
\end{learningcard}

\begin{learningcard}
\SoruHeader
Otoenkodör türlerinden hangisi gürültülü veriden temiz veriyi öğrenmek için tasarlanmıştır?
\par\vspace{\baselineskip} \tcblower
\CevapHeader
Denoising Autoencoder.
\end{learningcard}

\begin{learningcard}
\SoruHeader
Otoenkodörlerde (Autoencoders) 'Bottleneck' katmanının görevi nedir?
\par\vspace{\baselineskip} \tcblower
\CevapHeader
Veriyi en düşük boyutlu (sıkıştırılmış) temsiline ($z$) indirgemek.
\end{learningcard}

\begin{learningcard}
\SoruHeader
Otoenkodörlerde (Autoencoders) verinin en düşük boyutlu temsilinin bulunduğu katmana ne denir?
\par\vspace{\baselineskip} \tcblower
\CevapHeader
Bottleneck (Darboğaz).
\end{learningcard}

\begin{learningcard}
\SoruHeader
PCA (Temel Bileşen Analizi) algoritmasının ilk adımı nedir?
\par\vspace{\baselineskip} \tcblower
\CevapHeader
Veriyi merkezlemek (sütun ortalamasını her gözlemden çıkarmak).
\end{learningcard}

\begin{learningcard}
\SoruHeader
PCA algoritmasında özdeğerlerin ($\lambda_i$) büyüklüğü neyi temsil eder?
\par\vspace{\baselineskip} \tcblower
\CevapHeader
İlgili bileşenin taşıdığı varyans miktarını.
\end{learningcard}

\begin{learningcard}
\SoruHeader
PCA formülasyonunda özvektörler arasındaki ilişki nasıldır?
\par\vspace{\baselineskip} \tcblower
\CevapHeader
Özvektörler birbirine ortogonaldir (diktir) ve aralarında korelasyon yoktur.
\end{learningcard}

\begin{learningcard}
\SoruHeader
PCA'da açıklanan varyans oranı (EVR) nasıl hesaplanır?
\par\vspace{\baselineskip} \tcblower
\CevapHeader
$EVR(k) = \frac{\sum_{i=1}^k \lambda_i}{\sum_{i=1}^d \lambda_i}$
\end{learningcard}

\begin{learningcard}
\SoruHeader
PCA'da en uygun bileşen sayısını ($k$) belirlemek için kullanılan görsel grafiğe ne ad verilir?
\par\vspace{\baselineskip} \tcblower
\CevapHeader
Scree Plot (Yamaç Grafiği).
\end{learningcard}

\begin{learningcard}
\SoruHeader
PCA'da her bir ana bileşenin (PC) bir sonrakine göre geometrik konumu nedir?
\par\vspace{\baselineskip} \tcblower
\CevapHeader
Birbirlerine diktirler (Ortogonal).
\end{learningcard}

\begin{learningcard}
\SoruHeader
PCA'da kovaryans matrisi $C$ hangi matematiksel formülle hesaplanır?
\par\vspace{\baselineskip} \tcblower
\CevapHeader
$C = \frac{1}{N} \tilde{X}^T \tilde{X}$ (Burada $\tilde{X}$ merkezlenmiş veridir).
\end{learningcard}

\begin{learningcard}
\SoruHeader
PCA'da kullanılan 'Scree Plot' grafiği neyi belirlemek için kullanılır?
\par\vspace{\baselineskip} \tcblower
\CevapHeader
Özdeğerlerin sıralı dizilimindeki 'dirsek' noktasını bularak optimal bileşen sayısını ($k$) seçmek.
\end{learningcard}

\begin{learningcard}
\SoruHeader
PCA'da verinin merkezlenmesi için hangi işlem yapılır?
\par\vspace{\baselineskip} \tcblower
\CevapHeader
$\tilde{X} = X - \mu$ (Her özellikten sütun ortalamasının çıkarılması).
\end{learningcard}

\begin{learningcard}
\SoruHeader
PCA'nın geometrik yorumunda birinci temel bileşen (PC1) hangi yöndedir?
\par\vspace{\baselineskip} \tcblower
\CevapHeader
Verideki maksimum varyansın olduğu doğrultudadır.
\end{learningcard}

\begin{learningcard}
\SoruHeader
R-Kare ($R^2$) skoru neyi ifade eder?
\par\vspace{\baselineskip} \tcblower
\CevapHeader
Hedef değişkendeki varyansın yüzde kaçının modeldeki özellikler tarafından açıklandığını.
\end{learningcard}

\begin{learningcard}
\SoruHeader
Random Forest'ta her düğümde rastgele $m < d$ özellik seçilmesinin temel amacı nedir?
\par\vspace{\baselineskip} \tcblower
\CevapHeader
Ağaçlar arasındaki korelasyonu düşürerek topluluk varyansını azaltmak.
\end{learningcard}

\begin{learningcard}
\SoruHeader
Rastgele Orman (Random Forest) algoritmasındaki 'Bagging' mantığı nedir?
\par\vspace{\baselineskip} \tcblower
\CevapHeader
Eğitim verisinden tekrarlı rastgele örneklemler (bootstrap) alarak birden fazla ağaç eğitmek ve sonuçları oylamak.
\end{learningcard}

\begin{learningcard}
\SoruHeader
Rastgele Orman (Random Forest) algoritmasının temel çalışma prensibi nedir?
\par\vspace{\baselineskip} \tcblower
\CevapHeader
Birden fazla karar ağacını bootstrap örnekleme ve rastgele özellik seçimiyle birleştirmek.
\end{learningcard}

\begin{learningcard}
\SoruHeader
Regresyon modelinde $R^2 = 0.82$ sonucu ne anlama gelir?
\par\vspace{\baselineskip} \tcblower
\CevapHeader
Model, verideki değişkenliğin $\%82$'sini açıklayabilmektedir.
\end{learningcard}

\begin{learningcard}
\SoruHeader
Regresyon modellerinde 'Artık Değer' (Residual) nedir?
\par\vspace{\baselineskip} \tcblower
\CevapHeader
Gerçek değer ($y_i$) ile tahmin edilen değer ($\hat{y}_i$) arasındaki fark.
\end{learningcard}

\begin{learningcard}
\SoruHeader
RMSE metriği MSE'ye göre neden daha yorumlanabilirdir?
\par\vspace{\baselineskip} \tcblower
\CevapHeader
Karekök işlemi sayesinde hata birimini orijinal verinin birimine döndürdüğü için.
\end{learningcard}

\begin{learningcard}
\SoruHeader
Sigmoid fonksiyonunun türevinin çok küçük olması hangi eğitim problemine yol açar?
\par\vspace{\baselineskip} \tcblower
\CevapHeader
Gradyan Kaybolması (Vanishing Gradient).
\end{learningcard}

\begin{learningcard}
\SoruHeader
Sinir ağlarında 'Batch Normalization' neden kullanılır?
\par\vspace{\baselineskip} \tcblower
\CevapHeader
Eğitimi hızlandırmak ve gradyan sorunlarını (vanishing gradient) azaltmak için katman girişlerini normalize eder.
\end{learningcard}

\begin{learningcard}
\SoruHeader
Sinir ağlarında 'Dropout' tekniği neden kullanılır?
\par\vspace{\baselineskip} \tcblower
\CevapHeader
Nöronları rastgele kapatarak aşırı öğrenmeyi (overfitting) engellemek için.
\end{learningcard}

\begin{learningcard}
\SoruHeader
Sinir ağlarında 'ReLU' aktivasyon fonksiyonunun tanımı nedir?
\par\vspace{\baselineskip} \tcblower
\CevapHeader
$g(z) = \max(0, z)$.
\end{learningcard}

\begin{learningcard}
\SoruHeader
Sinir ağlarında 'Softmax' fonksiyonu genellikle hangi katmanda kullanılır?
\par\vspace{\baselineskip} \tcblower
\CevapHeader
Çok sınıflı sınıflandırma problemlerinde çıkış katmanında.
\end{learningcard}

\begin{learningcard}
\SoruHeader
Sinir ağlarında katmanlar boyunca verinin işlenerek tahmin üretilmesi sürecine ne ad verilir?
\par\vspace{\baselineskip} \tcblower
\CevapHeader
İleri Besleme (Forward Propagation).
\end{learningcard}

\begin{learningcard}
\SoruHeader
Sinir ağlarında negatif değerleri sıfıra eşitleyen en yaygın aktivasyon fonksiyonu hangisidir?
\par\vspace{\baselineskip} \tcblower
\CevapHeader
ReLU ($g(z) = \max(0, z)$).
\end{learningcard}

\begin{learningcard}
\SoruHeader
SOM (Kendi Kendini Organize Eden Haritalar) sinir ağı mimarisinin temel işlevi nedir?
\par\vspace{\baselineskip} \tcblower
\CevapHeader
Yüksek boyutlu veriyi topolojiyi koruyarak düşük boyutlu (genellikle 2D) bir haritaya indirgemek.
\end{learningcard}

\begin{learningcard}
\SoruHeader
SVD (Tekil Değer Ayrışımı) ile PCA arasındaki ilişki nedir?
\par\vspace{\baselineskip} \tcblower
\CevapHeader
SVD, PCA hesaplamasını sayısal olarak daha kararlı bir şekilde gerçekleştirmek için kullanılır.
\end{learningcard}

\begin{learningcard}
\SoruHeader
SVM algoritmasında 'Soft Margin' kullanımı neyi amaçlar?
\par\vspace{\baselineskip} \tcblower
\CevapHeader
Verideki bazı gürültülü noktaların yanlış sınıfta kalmasına izin vererek daha genel bir karar sınırı oluşturmayı.
\end{learningcard}

\begin{learningcard}
\SoruHeader
SVM'de 'Kernel Trick' ne işe yarar?
\par\vspace{\baselineskip} \tcblower
\CevapHeader
Veriyi daha yüksek boyutlu bir özellik uzayına taşıyarak doğrusal olmayan sınırların çizilmesini sağlar.
\end{learningcard}

\begin{learningcard}
\SoruHeader
SVM'de veriyi doğrusal olmayan bir üst uzaya taşıyarak ayrıştırmayı sağlayan yönteme ne denir?
\par\vspace{\baselineskip} \tcblower
\CevapHeader
Kernel Trick (Çekirdek Yöntemi).
\end{learningcard}

\begin{learningcard}
\SoruHeader
Sınıflandırma (Classification) problemi makine öğrenmesinin hangi kategorisine girer?
\par\vspace{\baselineskip} \tcblower
\CevapHeader
Denetimli Öğrenme (Supervised Learning).
\end{learningcard}

\begin{learningcard}
\SoruHeader
Sınıflandırma performansını ölçen Karmaşıklık Matrisi'nde (Confusion Matrix) FN neyi temsil eder?
\par\vspace{\baselineskip} \tcblower
\CevapHeader
Yanlış Negatif (Gerçekte pozitif olanın negatif tahmin edilmesi).
\end{learningcard}

\begin{learningcard}
\SoruHeader
Sınıflandırma probleminde modelin çıktı olarak ürettiği olasılığı 0 veya 1'e dönüştüren fonksiyona ne denir?
\par\vspace{\baselineskip} \tcblower
\CevapHeader
Sigmoid Fonksiyonu: $\sigma(z) = \frac{1}{1 + e^{-z}}$.
\end{learningcard}

\begin{learningcard}
\SoruHeader
t-SNE algoritması yüksek boyuttaki benzerlikleri düşük boyuta taşırken hangi olasılık dağılımını kullanır?
\par\vspace{\baselineskip} \tcblower
\CevapHeader
t-dağılımı (Student's t-distribution).
\end{learningcard}

\begin{learningcard}
\SoruHeader
UMAP algoritmasının t-SNE'ye göre en büyük pratik avantajı nedir?
\par\vspace{\baselineskip} \tcblower
\CevapHeader
Daha hızlı çalışması ve yeni veri noktalarının projeksiyonuna izin vermesi.
\end{learningcard}

\begin{learningcard}
\SoruHeader
UMAP yönteminin t-SNE'ye göre avantajı nedir?
\par\vspace{\baselineskip} \tcblower
\CevapHeader
Daha hızlı çalışması ve yeni veri noktalarını mevcut projeksiyona dahil edebilmesidir.
\end{learningcard}

\begin{learningcard}
\SoruHeader
Vektörleştirme (Vectorization) sinir ağı eğitiminde neden önemlidir?
\par\vspace{\baselineskip} \tcblower
\CevapHeader
Çok sayıda skaler işlemi tek bir matris çarpımına dönüştürerek donanım (GPU) hızlandırması sağlar.
\end{learningcard}

\begin{learningcard}
\SoruHeader
Veriyi 2D veya 3D'de görselleştirmek için t-SNE'nin kullandığı temel benzerlik ölçütü nedir?
\par\vspace{\baselineskip} \tcblower
\CevapHeader
Koşullu olasılıklar tabanlı benzerlik (t-dağılımı).
\end{learningcard}

\begin{learningcard}
\SoruHeader
XOR problemi neden tek katmanlı (doğrusal) modellerle çözülemez?
\par\vspace{\baselineskip} \tcblower
\CevapHeader
XOR verisi doğrusal olarak ayrıştırılamaz, modellemek için doğrusal olmayan bir gizli katman gerekir.
\end{learningcard}

\begin{learningcard}
\SoruHeader
Yapay sinir ağlarında 'Dropout' yönteminin işlevi nedir?
\par\vspace{\baselineskip} \tcblower
\CevapHeader
Eğitim sırasında nöronları rastgele kapatarak modelin belirli nöronlara aşırı bağımlı olmasını (overfitting) engellemek.
\end{learningcard}

\begin{learningcard}
\SoruHeader
Yapay sinir ağlarında bir nöronun temel matematiksel denklemi nedir?
\par\vspace{\baselineskip} \tcblower
\CevapHeader
$a = \phi(\sum w_j x_j + b)$ (Burada $\phi$ aktivasyon fonksiyonudur).
\end{learningcard}

\begin{learningcard}
\SoruHeader
Yapay sinir ağlarında tek bir nöronun matematiksel denklemi nedir?
\par\vspace{\baselineskip} \tcblower
\CevapHeader
$a = \phi(\sum w_j x_j + b)$
\end{learningcard}

\begin{learningcard}
\SoruHeader
Yüksek boyutlu veriyi 2D veya 3D uzayda görselleştirmek için en sık tercih edilen doğrusal olmayan yöntemler hangileridir?
\par\vspace{\baselineskip} \tcblower
\CevapHeader
t-SNE ve UMAP yöntemleri.
\end{learningcard}

\begin{learningcard}
\SoruHeader
Çok sınıflı sınıflandırma problemlerinde kullanılan hata fonksiyonu hangisidir?
\par\vspace{\baselineskip} \tcblower
\CevapHeader
Çapraz Entropi (Cross-Entropy Loss).
\end{learningcard}

\begin{learningcard}
\SoruHeader
Özelliklerin dönüştürülerek yeni yapay bileşenlerin üretildiği işleme ne ad verilir?
\par\vspace{\baselineskip} \tcblower
\CevapHeader
Boyut Düşürme (Dimensionality Reduction).
\end{learningcard}

\end{document}
